\chapter{世界はどう内側に入るか}
% full version: chapters/11_sensors.tex

どの感覚器も、マイクや光電セルのような変換器である。光、音、あるいは分子がセンサー細胞の特別なタンパク質に働きかけ、タンパク質が「蛇口」を開き、電流が流れ、連鎖の終わりでカチッが生まれてネットワークへ送られる。ほぼすべてのセンサーはグルタミン酸を放出するので、マシンの入力はそのまま電話網につながっている。

\section*{物理の限界で}

脳のセンサーは、可能なことの限界ぎりぎりで働いている。光のセンサーは、暗闇でただ1個の光の粒子に応える。音のセンサーは、1ナノメートル未満、つまり原子数個分より小さい変位に気づく。薄暗がりでは、精度を制限するのはもはや目ではなく光そのものだ。粒子はランダムに届くので、色合いを見分けるにはたくさんためる必要がある。だから夕暮れには見るのが遅くなる。目は、長時間露光のカメラのように、長く光をためるのだ。

\section*{自動露出}

明るさは星明かりの夜から日に照らされた雪まで数十億倍も変わり、音の大きさはそれ以上に変わる。一方、カチッの頻度は毎秒0から数百回までしか変わらない。一方を他方にそのまま押し込むのは不可能だ。解決策はカメラと同じ、自動露出である。センサーは周囲の平均レベルに合わせて感度を絶えず調整し、明るさそのものではなく\emph{背景と比べた}明るさを伝える。一定のものはやがて応答を引き起こさなくなり、変化はそのたびに応答を引き起こす。マシンの入力は、最初から変化について語っているのだ。

\pencilfig{125}{%
% light rises in two tenfold steps, then drops a hundredfold; sensor rate = baseline + gain * (log light - adapted level),
% the adapted level follows log light with a 0.7 s time constant; rate clipped at zero
\draw[penlite=pcAmber] (0,1.75) -- (1.4,1.75) -- (1.4,2.1) -- (4.2,2.1) -- (4.2,2.45) -- (6.3,2.45) -- (6.3,1.75) -- (8.4,1.75);
\node[lbl, text=pcAmber!70!black, anchor=east] at (-0.05,2.0) {光};
\node[lbl, anchor=south] at (2.8,2.12) {10倍明るく}; \node[lbl, anchor=south] at (5.25,2.47) {さらに10倍};
\node[lbl, anchor=south] at (7.35,1.77) {再び暗く};
\draw[pen] (0,0) -- (8.4,0);
\draw[pcGray, densely dashed, line width=0.5pt] (0,0.32) -- (8.4,0.32);
\draw[penlite=pcBlue] plot coordinates {(0.00,0.32) (1.26,0.32) (1.40,1.02) (1.54,0.85) (1.68,0.72) (1.82,0.62) (1.96,0.54) (2.10,0.49) (2.24,0.45) (2.38,0.41) (2.52,0.39) (2.66,0.37) (2.80,0.36) (3.08,0.34) (3.50,0.33) (4.06,0.32) (4.20,1.03) (4.34,0.85) (4.48,0.72) (4.62,0.62) (4.76,0.54) (4.90,0.49) (5.04,0.45) (5.18,0.41) (5.32,0.39) (5.46,0.37) (5.60,0.36) (5.88,0.34) (6.16,0.33) (6.30,0.00) (7.00,0.00) (7.14,0.07) (7.28,0.13) (7.42,0.18) (7.56,0.22) (7.70,0.24) (7.84,0.26) (7.98,0.28) (8.12,0.29) (8.40,0.30)};
\node[lbl, text=pcBlue, anchor=east] at (-0.05,0.32) {応答};
\node[lbl, anchor=north] at (6.65,-0.02) {沈黙};
}{光は1段ごとに10倍ずつ明るくなる。センサーは変化のたびに急上昇で応え、すぐに元のレベルに戻る。再び暗くなると、しばらく黙り込む。}

エンジニアはこの手法をイベントカメラで再現した。このカメラは時計に合わせてフレームを撮るのではない。各画素が、何かが変わったときに自分で「明るくなった」「暗くなった」と伝える。第2章の「はい」と「いいえ」の配線の組そのものだが、シリコンでできている。

\section*{目は画像を圧縮する}

目には約1億個の光センサーがあるが、画像を脳へ運ぶ配線は約100万本しかない。画像は目を出る前に、およそ100分の1に圧縮される。平坦な部分はほとんどコストがかからず、伝えられるのは輪郭と変化だ。動物の目で行われ、ヒトに換算された見積もりによれば、目は脳へ毎秒約1000万ビットを送る。昔のネットワークケーブル並みである。

\section*{耳はピアノ}

耳は、エンジニアならフィルタバンクを置くように、音を周波数に分解する。耳の中の弾力のある仕切りは場所ごとに異なる周波数に共鳴し、各部分のセンサーは自分の「鍵盤」を聞いている。鳴ったセンサーの番号が音の高さである。

\pencilfig{11}{\pencilart{11}}{耳は逆向きのピアノのようにできている。各鍵盤が自分の音を聞いている。}

さらに、耳は受け身ではない。耳の細胞の一部は、音に合わせて自分で仕切りを揺らし、小さな音を数千倍に増幅するが、大きな音はほとんど増幅しない。これは自励発振の境目にある増幅器で、そこでは特に感度が高い。しかも低い周波数では、ニューロンが波に合わせてカチッと鳴り、音の正確なタイミングを保つ。脳はそれを使って音源の方向を見つける。周波数が高くなるほどニューロンはついていけなくなり、残るのは場所による符号だけになる。

\section*{ほかの感覚}

ほぼどの感覚にも、これまでの章の手法が見てとれる。温度は、温かさ用と冷たさ用の2本の別々の配線で伝えられる。触覚は、触れた瞬間に応答してすぐ黙るセンサーと、圧力がある限り応答を保つセンサーを使う。そして嗅覚はバーコードのようにできている。ヒトの嗅覚センサーの種類は約400で、においとは反応した種類の組み合わせである。そうした組み合わせの数は天文学的だ。

\fullversion{11}
